\section{评测协议设计：避免 ``高分低用''}

讲者批评了只报 pass rate 的评估方式。本节将该观点展开为可执行评测协议：在同一 test-time budget 下，同时汇报质量、成本、时延、稳定性与合规性，避免 ``靠更多计算换高分'' 的比较失真。

\begin{importantbox}{评测协议的第一原则}
任何跨模型比较都必须先对齐预算：包括最大 token、最大工具调用次数、最大运行时长。预算不对齐，分数不可比。
\end{importantbox}

\begin{table}[H]
\centering
\begin{tabular}{p{2.3cm}p{3.0cm}p{5.0cm}}
\toprule
\textbf{指标族} & \textbf{代表指标} & \textbf{解释与注意事项} \\
\midrule
任务质量 & pass@1, pass@k, solve rate & 必须注明预算与任务分布；建议同时报告中位数与尾部表现 \\
成本效率 & avg tokens, tool calls, GPU-seconds & 与质量成对展示，避免只优化一侧 \\
交互体验 & latency P50/P95, progress update rate & 对 agent 产品极其关键，影响用户留存 \\
鲁棒合规 & cheat rate, policy violation rate & 反映是否 ``正确地做对''，而非 ``侥幸做对'' \\
稳定性 & retry rate, infra failure rate & 排除系统噪声后再解读模型收益 \\
\bottomrule
\end{tabular}
\caption{建议的 Agent 评测指标矩阵}
\end{table}

我们可以把综合得分写成分层函数：
\[
S = w_q \cdot Q - w_c \cdot C - w_t \cdot T - w_r \cdot R
\]
其中 \(Q\) 是任务质量，\(C\) 是成本，\(T\) 是时延，\(R\) 是风险（违规与作弊率）。该函数不是为了追求绝对数学正确，而是强制团队在汇报中显式说明 trade-off。

\begin{knowledgebox}{评测汇报模板（推荐）}
\begin{itemize}
    \item 固定预算：每题最多 20k tokens、最多 20 次工具调用、最长 10 分钟；
    \item 主指标：pass@1、pass@k、P95 latency、cheat rate；
    \item 辅指标：平均 token 成本、平均工具调用数、失败回放样例；
    \item 对比维度：与上一个线上版本做等预算 side-by-side。
\end{itemize}
\end{knowledgebox}

此外，评测集需分层抽样：简单题、中等题、长链任务、对抗样本、工具未知样本。只在 ``干净样本'' 上评估会显著高估模型上线表现。

\begin{warningbox}{最常见的评测偏差}
将 ``模型更长思考带来的提升'' 误判为 ``模型能力提升''。如果测试预算翻倍，pass@k 上升并不一定意味着训练策略更优。
\end{warningbox}

\subsection{本章小结}
高质量评测协议必须预算对齐、指标成组、样本分层。这样才能区分 ``真实能力进步'' 与 ``预算膨胀幻觉''。

